% --- Archon physics formalization source begin ---
% archon:physics
% archon:covers PhyXMiniProblems/problem_phyx_mini_0703.lean
% archon:source-report reports/phyx_mini/problem_phyx_mini_0703.source.json

\chapter{Physics problem phyx\_mini\_0703}
\label{ch:PhyXMiniProblems_problem_phyx_mini_0703}

\paragraph{Problem source.}
\#\# Physical scenario

Far out in space, a \textbackslash{}(100\{,\}000\textbackslash{},\textbackslash{}mathrm\{kg\}\textbackslash{}) rocket and a \textbackslash{}(200\{,\}000\textbackslash{},\textbackslash{}mathrm\{kg\}\textbackslash{}) rocket are docked at opposite ends of a motionless \textbackslash{}(90\textbackslash{},\textbackslash{}mathrm\{m\}\textbackslash{})-long connecting tunnel. The tunnel is rigid and its mass is much less than that of either rocket. The rockets start their engines simultaneously, each generating \textbackslash{}(50\{,\}000\textbackslash{},\textbackslash{}mathrm\{N\}\textbackslash{}) of thrust in opposite directions.

\#\# Question

What is the structure's angular velocity after \textbackslash{}(30\textbackslash{},\textbackslash{}mathrm\{s\}\textbackslash{})?

\#\# Answer choices

- A: A: \$0.00850\textbackslash{}

- B: B: \$0.00800\textbackslash{}

- C: C: \$0.00833\textbackslash{}

- D: D: \$0.00733\textbackslash{}

\#\# Figure caption (auxiliary; use the image as primary evidence)

The image is a diagram illustrating a physics problem involving two rockets and a tunnel.

1. **Objects:**
   - Two rockets, labeled with forces and masses.
   - A tunnel spanning 90 meters.
   - A reference line labeled \textbackslash{}(x\textbackslash{}) with specific positions marked.

2. **Rockets:**
   - The left rocket has a mass \textbackslash{}(m\_1 = 100,000 \textbackslash{}, \textbackslash{}text\{kg\}\textbackslash{}) and a force \textbackslash{}(\textbackslash{}vec\{F\}\_1 = 50,000 \textbackslash{}, \textbackslash{}text\{N\}\textbackslash{}) acting downward.
   - The right rocket has a mass \textbackslash{}(m\_2 = 200,000 \textbackslash{}, \textbackslash{}text\{kg\}\textbackslash{}) and a force \textbackslash{}(\textbackslash{}vec\{F\}\_2 = 50,000 \textbackslash{}, \textbackslash{}text\{N\}\textbackslash{}) acting upward.

3. **Positioning:**
   - Left rocket position \textbackslash{}(x\_1 = 0\textbackslash{}).
   - Right rocket position \textbackslash{}(x\_2 = 90 \textbackslash{}, \textbackslash{}text\{m\}\textbackslash{}).
   - Center of mass labeled as \textbackslash{}(x\_\textbackslash{}text\{cm\}\textbackslash{}).

4. **Additional Labels:**
   - Distances \textbackslash{}(r\_1\textbackslash{}) and \textbackslash{}(r\_2\textbackslash{}) are marked from the center of mass to each rocket.
   - A curved arrow near each rocket indicates exhaust or propulsion.

5. **Text:**
   - Labels and measurements for forces, masses, and positions provide details for solving a physics problem related to motion or forces.

\#\# Dataset metadata

Rotational Motion; Physical Model Grounding Reasoning, Multi-Formula Reasoning

\paragraph{Recorded answer/context.}
C: C: \$0.00833\textbackslash{}

\paragraph{Figure/image path.}
/root/proposal\_for\_physic/science-mango/phyx\_mini\_run/phyx\_data/test\_image/703.png

\paragraph{Formalization target.}
create a compiling Lean file with sorry bodies at `PhyXMiniProblems/problem\_phyx\_mini\_0703.lean`. The Lean declarations must preserve the physical quantities, dimensions or dimensional roles, figure labels, governing-law hypotheses, and final relation expressed by this problem.
Use Mathlib/Physlib names found through LeanExplore where available. If a physics API is missing, introduce faithful local abstractions rather than scalar placeholder aliases.

\begin{theorem}[Physics formalization target]
\label{thm:physics:phyx_mini_0703:target}
\lean{PhyXMiniProblems.ProblemPhyXMini0703.constantThrustPrediction_angularVelocity_after_thirty_seconds}
\uses{def:physics:phyx-mini-0703:phyxminiproblems-problemphyxmini0703-angularvelocityinradianspersecond, def:physics:phyx-mini-0703:phyxminiproblems-problemphyxmini0703-dockedrocketssetup, def:physics:phyx-mini-0703:phyxminiproblems-problemphyxmini0703-matchesdockedrocketsscenario, def:physics:phyx-mini-0703:phyxminiproblems-problemphyxmini0703-matchesdockedrocketsfigureanddata, def:physics:phyx-mini-0703:phyxminiproblems-problemphyxmini0703-hasphysicaldockedrocketsparameters, def:physics:phyx-mini-0703:phyxminiproblems-problemphyxmini0703-satisfiestworocketcenterofmassgeometry, def:physics:phyx-mini-0703:phyxminiproblems-problemphyxmini0703-satisfiesconstantthrustrotationaldynamics}
The assigned autoformalize agent should translate this physics problem into Lean declarations in the covered file, with theorem and lemma proof bodies written as `by sorry`.
\end{theorem}
\begin{proof}
This is an autoformalization task, not a proof task. Produce statements that can later be proved by the physics prover without weakening the physical model.
\end{proof}

% --- Archon declaration topology begin ---
\section{Declaration topology}
\begin{definition}[Planar Vector]
  \label{def:physics:phyx-mini-0703:phyxminiproblems-problemphyxmini0703-planarvector}
  \lean{PhyXMiniProblems.ProblemPhyXMini0703.PlanarVector}
  The two-dimensional Euclidean vector space of the supplied diagram
\end{definition}
\begin{proof}
  This is the declared model definition of Planar Vector.
\end{proof}
\begin{definition}[Mass Quantity]
  \label{def:physics:phyx-mini-0703:phyxminiproblems-problemphyxmini0703-massquantity}
  \lean{PhyXMiniProblems.ProblemPhyXMini0703.MassQuantity}
  A nonnegative, unit-independent physical mass
\end{definition}
\begin{proof}
  This is the declared model definition of Mass Quantity.
\end{proof}
\begin{definition}[Length Quantity]
  \label{def:physics:phyx-mini-0703:phyxminiproblems-problemphyxmini0703-lengthquantity}
  \lean{PhyXMiniProblems.ProblemPhyXMini0703.LengthQuantity}
  A nonnegative, unit-independent physical length
\end{definition}
\begin{proof}
  This is the declared model definition of Length Quantity.
\end{proof}
\begin{definition}[Time Quantity]
  \label{def:physics:phyx-mini-0703:phyxminiproblems-problemphyxmini0703-timequantity}
  \lean{PhyXMiniProblems.ProblemPhyXMini0703.TimeQuantity}
  A nonnegative, unit-independent time coordinate or elapsed duration
\end{definition}
\begin{proof}
  This is the declared model definition of Time Quantity.
\end{proof}
\begin{definition}[force Dimension]
  \label{def:physics:phyx-mini-0703:phyxminiproblems-problemphyxmini0703-forcedimension}
  \lean{PhyXMiniProblems.ProblemPhyXMini0703.forceDimension}
  The physical dimension of force, mass * length / time²
\end{definition}
\begin{proof}
  This is the declared model definition of force Dimension.
\end{proof}
\begin{definition}[Force Vector Quantity]
  \label{def:physics:phyx-mini-0703:phyxminiproblems-problemphyxmini0703-forcevectorquantity}
  \lean{PhyXMiniProblems.ProblemPhyXMini0703.ForceVectorQuantity}
  \uses{def:physics:phyx-mini-0703:phyxminiproblems-problemphyxmini0703-planarvector, def:physics:phyx-mini-0703:phyxminiproblems-problemphyxmini0703-forcedimension}
  A unit-independent planar force vector
\end{definition}
\begin{proof}
  This is the declared model definition of Force Vector Quantity.
\end{proof}
\begin{definition}[moment Of Inertia Dimension]
  \label{def:physics:phyx-mini-0703:phyxminiproblems-problemphyxmini0703-momentofinertiadimension}
  \lean{PhyXMiniProblems.ProblemPhyXMini0703.momentOfInertiaDimension}
  The physical dimension of a moment of inertia, mass * length²
\end{definition}
\begin{proof}
  This is the declared model definition of moment Of Inertia Dimension.
\end{proof}
\begin{definition}[Moment Of Inertia Quantity]
  \label{def:physics:phyx-mini-0703:phyxminiproblems-problemphyxmini0703-momentofinertiaquantity}
  \lean{PhyXMiniProblems.ProblemPhyXMini0703.MomentOfInertiaQuantity}
  \uses{def:physics:phyx-mini-0703:phyxminiproblems-problemphyxmini0703-momentofinertiadimension}
  A nonnegative, unit-independent moment of inertia
\end{definition}
\begin{proof}
  This is the declared model definition of Moment Of Inertia Quantity.
\end{proof}
\begin{definition}[torque Dimension]
  \label{def:physics:phyx-mini-0703:phyxminiproblems-problemphyxmini0703-torquedimension}
  \lean{PhyXMiniProblems.ProblemPhyXMini0703.torqueDimension}
  The physical dimension of torque, mass * length² / time²
\end{definition}
\begin{proof}
  This is the declared model definition of torque Dimension.
\end{proof}
\begin{definition}[Signed Torque Quantity]
  \label{def:physics:phyx-mini-0703:phyxminiproblems-problemphyxmini0703-signedtorquequantity}
  \lean{PhyXMiniProblems.ProblemPhyXMini0703.SignedTorqueQuantity}
  \uses{def:physics:phyx-mini-0703:phyxminiproblems-problemphyxmini0703-torquedimension}
  A signed torque about the axis normal to the diagram
\end{definition}
\begin{proof}
  This is the declared model definition of Signed Torque Quantity.
\end{proof}
\begin{definition}[Angular Velocity Quantity]
  \label{def:physics:phyx-mini-0703:phyxminiproblems-problemphyxmini0703-angularvelocityquantity}
  \lean{PhyXMiniProblems.ProblemPhyXMini0703.AngularVelocityQuantity}
  Signed angular velocity; radians are dimensionless
\end{definition}
\begin{proof}
  This is the declared model definition of Angular Velocity Quantity.
\end{proof}
\begin{definition}[Angular Acceleration Quantity]
  \label{def:physics:phyx-mini-0703:phyxminiproblems-problemphyxmini0703-angularaccelerationquantity}
  \lean{PhyXMiniProblems.ProblemPhyXMini0703.AngularAccelerationQuantity}
  Signed angular acceleration; radians are dimensionless
\end{definition}
\begin{proof}
  This is the declared model definition of Angular Acceleration Quantity.
\end{proof}
\begin{definition}[Velocity Vector Quantity]
  \label{def:physics:phyx-mini-0703:phyxminiproblems-problemphyxmini0703-velocityvectorquantity}
  \lean{PhyXMiniProblems.ProblemPhyXMini0703.VelocityVectorQuantity}
  \uses{def:physics:phyx-mini-0703:phyxminiproblems-problemphyxmini0703-planarvector}
  Unit-independent center-of-mass velocity in the diagram plane
\end{definition}
\begin{proof}
  This is the declared model definition of Velocity Vector Quantity.
\end{proof}
\begin{definition}[nonnegative SIReadout]
  \label{def:physics:phyx-mini-0703:phyxminiproblems-problemphyxmini0703-nonnegativesireadout}
  \lean{PhyXMiniProblems.ProblemPhyXMini0703.nonnegativeSIReadout}
  Coherent-SI readout of a nonnegative dimensionful scalar
\end{definition}
\begin{proof}
  This is the declared model definition of nonnegative SIReadout.
\end{proof}
\begin{definition}[signed SIReadout]
  \label{def:physics:phyx-mini-0703:phyxminiproblems-problemphyxmini0703-signedsireadout}
  \lean{PhyXMiniProblems.ProblemPhyXMini0703.signedSIReadout}
  Coherent-SI readout of a signed dimensionful scalar
\end{definition}
\begin{proof}
  This is the declared model definition of signed SIReadout.
\end{proof}
\begin{definition}[vector SIReadout]
  \label{def:physics:phyx-mini-0703:phyxminiproblems-problemphyxmini0703-vectorsireadout}
  \lean{PhyXMiniProblems.ProblemPhyXMini0703.vectorSIReadout}
  \uses{def:physics:phyx-mini-0703:phyxminiproblems-problemphyxmini0703-planarvector}
  Coherent-SI readout of a dimensionful planar vector
\end{definition}
\begin{proof}
  This is the declared model definition of vector SIReadout.
\end{proof}
\begin{definition}[mass In Kilograms]
  \label{def:physics:phyx-mini-0703:phyxminiproblems-problemphyxmini0703-massinkilograms}
  \lean{PhyXMiniProblems.ProblemPhyXMini0703.massInKilograms}
  \uses{def:physics:phyx-mini-0703:phyxminiproblems-problemphyxmini0703-massquantity, def:physics:phyx-mini-0703:phyxminiproblems-problemphyxmini0703-nonnegativesireadout}
  Mass readout in kilograms
\end{definition}
\begin{proof}
  This is the declared model definition of mass In Kilograms.
\end{proof}
\begin{definition}[length In Meters]
  \label{def:physics:phyx-mini-0703:phyxminiproblems-problemphyxmini0703-lengthinmeters}
  \lean{PhyXMiniProblems.ProblemPhyXMini0703.lengthInMeters}
  \uses{def:physics:phyx-mini-0703:phyxminiproblems-problemphyxmini0703-lengthquantity, def:physics:phyx-mini-0703:phyxminiproblems-problemphyxmini0703-nonnegativesireadout}
  Length or longitudinal-position readout in metres
\end{definition}
\begin{proof}
  This is the declared model definition of length In Meters.
\end{proof}
\begin{definition}[time In Seconds]
  \label{def:physics:phyx-mini-0703:phyxminiproblems-problemphyxmini0703-timeinseconds}
  \lean{PhyXMiniProblems.ProblemPhyXMini0703.timeInSeconds}
  \uses{def:physics:phyx-mini-0703:phyxminiproblems-problemphyxmini0703-timequantity, def:physics:phyx-mini-0703:phyxminiproblems-problemphyxmini0703-nonnegativesireadout}
  Time readout in seconds
\end{definition}
\begin{proof}
  This is the declared model definition of time In Seconds.
\end{proof}
\begin{definition}[inertia In Kilogram Meters Squared]
  \label{def:physics:phyx-mini-0703:phyxminiproblems-problemphyxmini0703-inertiainkilogrammeterssquared}
  \lean{PhyXMiniProblems.ProblemPhyXMini0703.inertiaInKilogramMetersSquared}
  \uses{def:physics:phyx-mini-0703:phyxminiproblems-problemphyxmini0703-momentofinertiaquantity, def:physics:phyx-mini-0703:phyxminiproblems-problemphyxmini0703-nonnegativesireadout}
  Moment-of-inertia readout in kilogram metres squared
\end{definition}
\begin{proof}
  This is the declared model definition of inertia In Kilogram Meters Squared.
\end{proof}
\begin{definition}[torque In Newton Meters]
  \label{def:physics:phyx-mini-0703:phyxminiproblems-problemphyxmini0703-torqueinnewtonmeters}
  \lean{PhyXMiniProblems.ProblemPhyXMini0703.torqueInNewtonMeters}
  \uses{def:physics:phyx-mini-0703:phyxminiproblems-problemphyxmini0703-signedtorquequantity, def:physics:phyx-mini-0703:phyxminiproblems-problemphyxmini0703-signedsireadout}
  Signed torque readout in newton metres
\end{definition}
\begin{proof}
  This is the declared model definition of torque In Newton Meters.
\end{proof}
\begin{definition}[angular Velocity In Radians Per Second]
  \label{def:physics:phyx-mini-0703:phyxminiproblems-problemphyxmini0703-angularvelocityinradianspersecond}
  \lean{PhyXMiniProblems.ProblemPhyXMini0703.angularVelocityInRadiansPerSecond}
  \uses{def:physics:phyx-mini-0703:phyxminiproblems-problemphyxmini0703-angularvelocityquantity, def:physics:phyx-mini-0703:phyxminiproblems-problemphyxmini0703-signedsireadout}
  Signed angular-velocity readout in radians per second
\end{definition}
\begin{proof}
  This is the declared model definition of angular Velocity In Radians Per Second.
\end{proof}
\begin{definition}[angular Acceleration In Radians Per Second Squared]
  \label{def:physics:phyx-mini-0703:phyxminiproblems-problemphyxmini0703-angularaccelerationinradianspersecondsquared}
  \lean{PhyXMiniProblems.ProblemPhyXMini0703.angularAccelerationInRadiansPerSecondSquared}
  \uses{def:physics:phyx-mini-0703:phyxminiproblems-problemphyxmini0703-angularaccelerationquantity, def:physics:phyx-mini-0703:phyxminiproblems-problemphyxmini0703-signedsireadout}
  Signed angular-acceleration readout in radians per second squared
\end{definition}
\begin{proof}
  This is the declared model definition of angular Acceleration In Radians Per Second Squared.
\end{proof}
\begin{definition}[longitudinal Axis]
  \label{def:physics:phyx-mini-0703:phyxminiproblems-problemphyxmini0703-longitudinalaxis}
  \lean{PhyXMiniProblems.ProblemPhyXMini0703.longitudinalAxis}
  \uses{def:physics:phyx-mini-0703:phyxminiproblems-problemphyxmini0703-planarvector}
  Unit vector along the tunnel's labelled longitudinal x axis
\end{definition}
\begin{proof}
  This is the declared model definition of longitudinal Axis.
\end{proof}
\begin{definition}[positive Transverse Axis]
  \label{def:physics:phyx-mini-0703:phyxminiproblems-problemphyxmini0703-positivetransverseaxis}
  \lean{PhyXMiniProblems.ProblemPhyXMini0703.positiveTransverseAxis}
  \uses{def:physics:phyx-mini-0703:phyxminiproblems-problemphyxmini0703-planarvector}
  Unit vector in the positive transverse direction in the diagram plane
\end{definition}
\begin{proof}
  This is the declared model definition of positive Transverse Axis.
\end{proof}
\begin{definition}[planar Cross Z]
  \label{def:physics:phyx-mini-0703:phyxminiproblems-problemphyxmini0703-planarcrossz}
  \lean{PhyXMiniProblems.ProblemPhyXMini0703.planarCrossZ}
  \uses{def:physics:phyx-mini-0703:phyxminiproblems-problemphyxmini0703-planarvector}
  The signed normal component of the planar cross product
\end{definition}
\begin{proof}
  This is the declared model definition of planar Cross Z.
\end{proof}
\begin{definition}[Rocket]
  \label{def:physics:phyx-mini-0703:phyxminiproblems-problemphyxmini0703-rocket}
  \lean{PhyXMiniProblems.ProblemPhyXMini0703.Rocket}
  The two rockets labelled 1 and 2 in the figure
\end{definition}
\begin{proof}
  This is the declared model definition of Rocket.
\end{proof}
\begin{definition}[Tunnel End]
  \label{def:physics:phyx-mini-0703:phyxminiproblems-problemphyxmini0703-tunnelend}
  \lean{PhyXMiniProblems.ProblemPhyXMini0703.TunnelEnd}
  The two physical ends of the connecting tunnel
\end{definition}
\begin{proof}
  This is the declared model definition of Tunnel End.
\end{proof}
\begin{definition}[Rotation Sense]
  \label{def:physics:phyx-mini-0703:phyxminiproblems-problemphyxmini0703-rotationsense}
  \lean{PhyXMiniProblems.ProblemPhyXMini0703.RotationSense}
  Sense of rotation about the normal axis of the supplied diagram
\end{definition}
\begin{proof}
  This is the declared model definition of Rotation Sense.
\end{proof}
\begin{definition}[Tunnel Model]
  \label{def:physics:phyx-mini-0703:phyxminiproblems-problemphyxmini0703-tunnelmodel}
  \lean{PhyXMiniProblems.ProblemPhyXMini0703.TunnelModel}
  Mechanical treatment assigned to the connecting tunnel
\end{definition}
\begin{proof}
  This is the declared model definition of Tunnel Model.
\end{proof}
\begin{definition}[Rocket Mass Model]
  \label{def:physics:phyx-mini-0703:phyxminiproblems-problemphyxmini0703-rocketmassmodel}
  \lean{PhyXMiniProblems.ProblemPhyXMini0703.RocketMassModel}
  Spatial treatment assigned to each rocket in the rotational model
\end{definition}
\begin{proof}
  This is the declared model definition of Rocket Mass Model.
\end{proof}
\begin{definition}[Engine Thrust Model]
  \label{def:physics:phyx-mini-0703:phyxminiproblems-problemphyxmini0703-enginethrustmodel}
  \lean{PhyXMiniProblems.ProblemPhyXMini0703.EngineThrustModel}
  Time dependence assigned to the thrust during the observation interval
\end{definition}
\begin{proof}
  This is the declared model definition of Engine Thrust Model.
\end{proof}
\begin{definition}[Docked Rockets Figure]
  \label{def:physics:phyx-mini-0703:phyxminiproblems-problemphyxmini0703-dockedrocketsfigure}
  \lean{PhyXMiniProblems.ProblemPhyXMini0703.DockedRocketsFigure}
  \uses{def:physics:phyx-mini-0703:phyxminiproblems-problemphyxmini0703-massquantity, def:physics:phyx-mini-0703:phyxminiproblems-problemphyxmini0703-lengthquantity, def:physics:phyx-mini-0703:phyxminiproblems-problemphyxmini0703-forcevectorquantity, def:physics:phyx-mini-0703:phyxminiproblems-problemphyxmini0703-rocket, def:physics:phyx-mini-0703:phyxminiproblems-problemphyxmini0703-rotationsense}
  Literal and qualitative content of image 703. The dimensional fields are the values printed beside the objects and arrows. The center-of-mass marker and the two lever-arm labels are visible, but the image prints no numerical values for them
\end{definition}
\begin{proof}
  This is the declared model definition of Docked Rockets Figure.
\end{proof}
\begin{definition}[Docked Rockets Setup]
  \label{def:physics:phyx-mini-0703:phyxminiproblems-problemphyxmini0703-dockedrocketssetup}
  \lean{PhyXMiniProblems.ProblemPhyXMini0703.DockedRocketsSetup}
  \uses{def:physics:phyx-mini-0703:phyxminiproblems-problemphyxmini0703-massquantity, def:physics:phyx-mini-0703:phyxminiproblems-problemphyxmini0703-lengthquantity, def:physics:phyx-mini-0703:phyxminiproblems-problemphyxmini0703-timequantity, def:physics:phyx-mini-0703:phyxminiproblems-problemphyxmini0703-forcevectorquantity, def:physics:phyx-mini-0703:phyxminiproblems-problemphyxmini0703-momentofinertiaquantity, def:physics:phyx-mini-0703:phyxminiproblems-problemphyxmini0703-signedtorquequantity, def:physics:phyx-mini-0703:phyxminiproblems-problemphyxmini0703-angularvelocityquantity, def:physics:phyx-mini-0703:phyxminiproblems-problemphyxmini0703-angularaccelerationquantity, def:physics:phyx-mini-0703:phyxminiproblems-problemphyxmini0703-velocityvectorquantity, def:physics:phyx-mini-0703:phyxminiproblems-problemphyxmini0703-rocket, def:physics:phyx-mini-0703:phyxminiproblems-problemphyxmini0703-tunnelend, def:physics:phyx-mini-0703:phyxminiproblems-problemphyxmini0703-tunnelmodel, def:physics:phyx-mini-0703:phyxminiproblems-problemphyxmini0703-rocketmassmodel, def:physics:phyx-mini-0703:phyxminiproblems-problemphyxmini0703-enginethrustmodel, def:physics:phyx-mini-0703:phyxminiproblems-problemphyxmini0703-dockedrocketsfigure}
  The fields record the physical assembly, its forces, the center-of-mass geometry, and its time-dependent motion. None is defined from a displayed answer. Relations among these independent fields are stated separately as geometry and dynamics laws below
\end{definition}
\begin{proof}
  This is the declared model definition of Docked Rockets Setup.
\end{proof}
\begin{definition}[Matches Docked Rockets Scenario]
  \label{def:physics:phyx-mini-0703:phyxminiproblems-problemphyxmini0703-matchesdockedrocketsscenario}
  \lean{PhyXMiniProblems.ProblemPhyXMini0703.MatchesDockedRocketsScenario}
  \uses{def:physics:phyx-mini-0703:phyxminiproblems-problemphyxmini0703-vectorsireadout, def:physics:phyx-mini-0703:phyxminiproblems-problemphyxmini0703-angularvelocityinradianspersecond, def:physics:phyx-mini-0703:phyxminiproblems-problemphyxmini0703-dockedrocketssetup}
  Qualitative physical setup stated in the prose
\end{definition}
\begin{proof}
  This is the declared model definition of Matches Docked Rockets Scenario.
\end{proof}
\begin{definition}[Matches Docked Rockets Figure And Data]
  \label{def:physics:phyx-mini-0703:phyxminiproblems-problemphyxmini0703-matchesdockedrocketsfigureanddata}
  \lean{PhyXMiniProblems.ProblemPhyXMini0703.MatchesDockedRocketsFigureAndData}
  \uses{def:physics:phyx-mini-0703:phyxminiproblems-problemphyxmini0703-vectorsireadout, def:physics:phyx-mini-0703:phyxminiproblems-problemphyxmini0703-massinkilograms, def:physics:phyx-mini-0703:phyxminiproblems-problemphyxmini0703-lengthinmeters, def:physics:phyx-mini-0703:phyxminiproblems-problemphyxmini0703-timeinseconds, def:physics:phyx-mini-0703:phyxminiproblems-problemphyxmini0703-positivetransverseaxis, def:physics:phyx-mini-0703:phyxminiproblems-problemphyxmini0703-dockedrocketssetup}
  Exact labels and arrow directions supplied by image 703 and the problem text. The transverse basis is chosen so rocket 2's arrow is positive. Consequently rocket 1's equal arrow is negative, exactly encoding the depicted opposite thrusts rather than only their magnitudes
\end{definition}
\begin{proof}
  This is the declared model definition of Matches Docked Rockets Figure And Data.
\end{proof}
\begin{definition}[Has Physical Docked Rockets Parameters]
  \label{def:physics:phyx-mini-0703:phyxminiproblems-problemphyxmini0703-hasphysicaldockedrocketsparameters}
  \lean{PhyXMiniProblems.ProblemPhyXMini0703.HasPhysicalDockedRocketsParameters}
  \uses{def:physics:phyx-mini-0703:phyxminiproblems-problemphyxmini0703-massinkilograms, def:physics:phyx-mini-0703:phyxminiproblems-problemphyxmini0703-lengthinmeters, def:physics:phyx-mini-0703:phyxminiproblems-problemphyxmini0703-timeinseconds, def:physics:phyx-mini-0703:phyxminiproblems-problemphyxmini0703-inertiainkilogrammeterssquared, def:physics:phyx-mini-0703:phyxminiproblems-problemphyxmini0703-dockedrocketssetup}
  Positivity and ordering assumptions selecting the physical configuration
\end{definition}
\begin{proof}
  This is the declared model definition of Has Physical Docked Rockets Parameters.
\end{proof}
\begin{definition}[Satisfies Two Rocket Center Of Mass Geometry]
  \label{def:physics:phyx-mini-0703:phyxminiproblems-problemphyxmini0703-satisfiestworocketcenterofmassgeometry}
  \lean{PhyXMiniProblems.ProblemPhyXMini0703.SatisfiesTwoRocketCenterOfMassGeometry}
  \uses{def:physics:phyx-mini-0703:phyxminiproblems-problemphyxmini0703-massinkilograms, def:physics:phyx-mini-0703:phyxminiproblems-problemphyxmini0703-lengthinmeters, def:physics:phyx-mini-0703:phyxminiproblems-problemphyxmini0703-dockedrocketssetup}
  Center-of-mass and endpoint geometry for two point masses on a massless rigid tunnel. The first law is the defining mass-moment balance, not the derived numerical center coordinate requested along the solution route
\end{definition}
\begin{proof}
  This is the declared model definition of Satisfies Two Rocket Center Of Mass Geometry.
\end{proof}
\begin{definition}[Satisfies Constant Thrust Rotational Dynamics]
  \label{def:physics:phyx-mini-0703:phyxminiproblems-problemphyxmini0703-satisfiesconstantthrustrotationaldynamics}
  \lean{PhyXMiniProblems.ProblemPhyXMini0703.SatisfiesConstantThrustRotationalDynamics}
  \uses{def:physics:phyx-mini-0703:phyxminiproblems-problemphyxmini0703-vectorsireadout, def:physics:phyx-mini-0703:phyxminiproblems-problemphyxmini0703-massinkilograms, def:physics:phyx-mini-0703:phyxminiproblems-problemphyxmini0703-lengthinmeters, def:physics:phyx-mini-0703:phyxminiproblems-problemphyxmini0703-timeinseconds, def:physics:phyx-mini-0703:phyxminiproblems-problemphyxmini0703-inertiainkilogrammeterssquared, def:physics:phyx-mini-0703:phyxminiproblems-problemphyxmini0703-torqueinnewtonmeters, def:physics:phyx-mini-0703:phyxminiproblems-problemphyxmini0703-angularvelocityinradianspersecond, def:physics:phyx-mini-0703:phyxminiproblems-problemphyxmini0703-angularaccelerationinradianspersecondsquared, def:physics:phyx-mini-0703:phyxminiproblems-problemphyxmini0703-longitudinalaxis, def:physics:phyx-mini-0703:phyxminiproblems-problemphyxmini0703-planarcrossz, def:physics:phyx-mini-0703:phyxminiproblems-problemphyxmini0703-dockedrocketssetup}
  Governing laws for the massless-tunnel, point-mass approximation: the net external force is the sum of the two thrust vectors; I = m₁ r₁² + m₂ r₂² about the center of mass; torque is the sum of the two signed planar lever-arm cross products; I α = τ; and constant angular acceleration integrates to ω(t) = ω(0) + α (t - 0) during the observation interval. No field states either the requested angular velocity or any answer-choice value
\end{definition}
\begin{proof}
  This is the declared model definition of Satisfies Constant Thrust Rotational Dynamics.
\end{proof}
\begin{lemma}[center Of Mass and lever Arm readouts]
  \label{lem:physics:phyx-mini-0703:phyxminiproblems-problemphyxmini0703-centerofmass-and-leverarm-readouts}
  \lean{PhyXMiniProblems.ProblemPhyXMini0703.centerOfMass_and_leverArm_readouts}
  \uses{def:physics:phyx-mini-0703:phyxminiproblems-problemphyxmini0703-lengthinmeters, def:physics:phyx-mini-0703:phyxminiproblems-problemphyxmini0703-dockedrocketssetup, def:physics:phyx-mini-0703:phyxminiproblems-problemphyxmini0703-matchesdockedrocketsfigureanddata, def:physics:phyx-mini-0703:phyxminiproblems-problemphyxmini0703-hasphysicaldockedrocketsparameters, def:physics:phyx-mini-0703:phyxminiproblems-problemphyxmini0703-satisfiestworocketcenterofmassgeometry}
  The unequal masses put the center of mass 60 m from rocket 1
\end{lemma}
\begin{proof}
  The conclusion follows from the governing relations and figure readouts named in the statement of center Of Mass and lever Arm readouts.
\end{proof}
\begin{lemma}[moment Of Inertia readout eq]
  \label{lem:physics:phyx-mini-0703:phyxminiproblems-problemphyxmini0703-momentofinertia-readout-eq}
  \lean{PhyXMiniProblems.ProblemPhyXMini0703.momentOfInertia_readout_eq}
  \uses{def:physics:phyx-mini-0703:phyxminiproblems-problemphyxmini0703-inertiainkilogrammeterssquared, def:physics:phyx-mini-0703:phyxminiproblems-problemphyxmini0703-dockedrocketssetup, def:physics:phyx-mini-0703:phyxminiproblems-problemphyxmini0703-matchesdockedrocketsfigureanddata, def:physics:phyx-mini-0703:phyxminiproblems-problemphyxmini0703-hasphysicaldockedrocketsparameters, def:physics:phyx-mini-0703:phyxminiproblems-problemphyxmini0703-satisfiestworocketcenterofmassgeometry, def:physics:phyx-mini-0703:phyxminiproblems-problemphyxmini0703-satisfiesconstantthrustrotationaldynamics}
  The point-mass moment of inertia is 5.40 × 10⁸ kg m²
\end{lemma}
\begin{proof}
  The conclusion follows from the governing relations and figure readouts named in the statement of moment Of Inertia readout eq.
\end{proof}
\begin{lemma}[net Torque readout eq]
  \label{lem:physics:phyx-mini-0703:phyxminiproblems-problemphyxmini0703-nettorque-readout-eq}
  \lean{PhyXMiniProblems.ProblemPhyXMini0703.netTorque_readout_eq}
  \uses{def:physics:phyx-mini-0703:phyxminiproblems-problemphyxmini0703-torqueinnewtonmeters, def:physics:phyx-mini-0703:phyxminiproblems-problemphyxmini0703-dockedrocketssetup, def:physics:phyx-mini-0703:phyxminiproblems-problemphyxmini0703-matchesdockedrocketsfigureanddata, def:physics:phyx-mini-0703:phyxminiproblems-problemphyxmini0703-hasphysicaldockedrocketsparameters, def:physics:phyx-mini-0703:phyxminiproblems-problemphyxmini0703-satisfiestworocketcenterofmassgeometry, def:physics:phyx-mini-0703:phyxminiproblems-problemphyxmini0703-satisfiesconstantthrustrotationaldynamics}
  The equal transverse thrusts form a 4.50 × 10⁶ N m torque couple
\end{lemma}
\begin{proof}
  The conclusion follows from the governing relations and figure readouts named in the statement of net Torque readout eq.
\end{proof}
\begin{lemma}[angular Acceleration readout eq]
  \label{lem:physics:phyx-mini-0703:phyxminiproblems-problemphyxmini0703-angularacceleration-readout-eq}
  \lean{PhyXMiniProblems.ProblemPhyXMini0703.angularAcceleration_readout_eq}
  \uses{def:physics:phyx-mini-0703:phyxminiproblems-problemphyxmini0703-angularaccelerationinradianspersecondsquared, def:physics:phyx-mini-0703:phyxminiproblems-problemphyxmini0703-dockedrocketssetup, def:physics:phyx-mini-0703:phyxminiproblems-problemphyxmini0703-matchesdockedrocketsfigureanddata, def:physics:phyx-mini-0703:phyxminiproblems-problemphyxmini0703-hasphysicaldockedrocketsparameters, def:physics:phyx-mini-0703:phyxminiproblems-problemphyxmini0703-satisfiestworocketcenterofmassgeometry, def:physics:phyx-mini-0703:phyxminiproblems-problemphyxmini0703-satisfiesconstantthrustrotationaldynamics}
  The governing laws give angular acceleration 1/120 rad/s²
\end{lemma}
\begin{proof}
  The conclusion follows from the governing relations and figure readouts named in the statement of angular Acceleration readout eq.
\end{proof}
\begin{definition}[Answer Choice]
  \label{def:physics:phyx-mini-0703:phyxminiproblems-problemphyxmini0703-answerchoice}
  \lean{PhyXMiniProblems.ProblemPhyXMini0703.AnswerChoice}
  Labels of the four displayed answer choices
\end{definition}
\begin{proof}
  This is the declared model definition of Answer Choice.
\end{proof}
\begin{definition}[displayed Angular Velocity Radians Per Second]
  \label{def:physics:phyx-mini-0703:phyxminiproblems-problemphyxmini0703-displayedangularvelocityradianspersecond}
  \lean{PhyXMiniProblems.ProblemPhyXMini0703.displayedAngularVelocityRadiansPerSecond}
  \uses{def:physics:phyx-mini-0703:phyxminiproblems-problemphyxmini0703-answerchoice}
  The decimal printed for each choice, interpreted in radians per second because the source question asks for angular velocity. These values are display data, not governing laws
\end{definition}
\begin{proof}
  This is the declared model definition of displayed Angular Velocity Radians Per Second.
\end{proof}
\begin{definition}[recorded Answer Choice]
  \label{def:physics:phyx-mini-0703:phyxminiproblems-problemphyxmini0703-recordedanswerchoice}
  \lean{PhyXMiniProblems.ProblemPhyXMini0703.recordedAnswerChoice}
  \uses{def:physics:phyx-mini-0703:phyxminiproblems-problemphyxmini0703-answerchoice}
  The answer label recorded by the dataset. It is source metadata only: no physical theorem identifies the observed angular velocity with this choice. Numerically, choice C is a three-decimal-place presentation of the angular acceleration 1/120 rad/s², whereas the source asks for angular velocity after 30 seconds. The physically supported answer to that question is the preceding theorem's 1/4 rad/s
\end{definition}
\begin{proof}
  This is the declared model definition of recorded Answer Choice.
\end{proof}
% --- Archon declaration topology end ---
% --- Archon physics formalization source end ---
